% Generated from locale/tr/content/normal-modal-logic/completeness/complete-consistent-sets.tex; do not hand-edit.


\olfileid[tr]{nml}{com}{ccs}

\olsection{Tam $\Sigma$-Tutarlı Kümeler}

$\Sigma$ bir modal !!{formula}s kümesi olsun; bunları bir normal modal
mantığın aksiyomları ya da tanımlayıcı ilkeleri olarak düşünün. Bir $\Gamma$
kümesi, ancak ve ancak $\Gamma\Proves/[\Sigma]\lfalse$ ise
$\Sigma$-tutarlıdır; yani her $!A_i\in\Gamma$ olmak üzere
$!A_1\lif(!A_2\lif\cdots(!A_n\lif\lfalse)\dots)$ formülünün $\Sigma$'dan
bir !!{derivation}{}i yoksa tutarlıdır. Her dünyanın özel türden bir
$\Sigma$-tutarlı küme olarak alındığı bir ``kanonik'' model kuracağız. Bu
kümeler yalnızca $\Sigma$-tutarlı değil, her modal !!{formula}{}ün doğruluk
değerini belirleyecek anlamda azami tutarlıdır: her $!A$ için ya
$!A\in\Gamma$ ya da $\lnot !A\in\Gamma$.

\begin{defn}
  Bir $\Gamma$ kümesi, ancak ve ancak $\Sigma$-tutarlıysa ve her $!A$ için
  ya $!A\in\Gamma$ ya da $\lnot !A\in\Gamma$ ise \emph{tam
  $\Sigma$-tutarlıdır}.
\end{defn}

Tam $\Sigma$-tutarlı $\Gamma$ kümelerinin bazı yararlı özellikleri vardır.
Öncelikle türetim altında kapalıdırlar: $\Gamma\Proves[\Sigma]!A$ ise
$!A\in\Gamma$. Özellikle bu, $!A\in\Sigma$ olan her !!a{formula}{}ün her
örneğinin de $\Gamma$ içinde olması demektir. Ayrıca $\Gamma$'ya üyelik,
önerme bağlaçlarının doğruluk koşullarını yansıtır. Bu, ``kanonik modeli''
tanımladığımızda önemli olacaktır.

\begin{prop}\ollabel{prop:ccs-properties}
  $\Gamma$ tam $\Sigma$-tutarlı olsun. O zaman:
  \begin{enumerate}
  \item \ollabel{prop:ccs-closed}%
    $\Gamma$, $\Sigma$ içinde türetim altında kapalıdır.
  \item \ollabel{prop:ccs-sigma}%
    $\Sigma\subseteq\Gamma$.
  \tagitem{prvFalse}{\ollabel{prop:ccs-lfalse}%
    $\lfalse\notin\Gamma$}{}
  \tagitem{prvTrue}{\ollabel{prop:ccs-ltrue}%
    $\ltrue\in\Gamma$}{}
  \item \ollabel{prop:ccs-lnot}%
    $\lnot!A\in\Gamma$ ancak ve ancak $!A\notin\Gamma$.
  \tagitem{prvAnd}{\ollabel{prop:ccs-land}%
    $!A\land !B\in\Gamma$ ancak ve ancak $!A\in\Gamma$ ve $!B\in\Gamma$}{}
  \tagitem{prvOr}{\ollabel{prop:ccs-lor}%
    $!A\lor !B\in\Gamma$ ancak ve ancak $!A\in\Gamma$ ya da $!B\in\Gamma$}{}
  \tagitem{prvIf}{\ollabel{prop:ccs-lif}%
    $!A\lif !B\in\Gamma$ ancak ve ancak $!A\notin\Gamma$ ya da $!B\in\Gamma$}{}
  \tagitem{prvIff}{\ollabel{prop:ccs-liff}%
    $!A\liff !B\in\Gamma$ ancak ve ancak ya $!A\in\Gamma$ ve
    $!B\in\Gamma$ ya da $!A\notin\Gamma$ ve $!B\notin\Gamma$}{}
  \end{enumerate}
\end{prop}

\begin{proof}
  \begin{enumerate}
  \item $\Gamma\Proves[\Sigma]!A$ fakat $!A\notin\Gamma$ olsun. $\Gamma$
    tam $\Sigma$-tutarlı olduğundan $\lnot!A\in\Gamma$ olur. Oysa
    $!A,\lnot !A\Proves[\Sigma]\lfalse$ olduğundan bu $\Gamma$'yı tutarsız
    yapardı.

  \item $!A\in\Sigma$ ise $\Gamma\Proves[\Sigma]!A$ ve türetim altında
    kapalılık, yani \olref{prop:ccs-closed} durumu gereğince $!A\in\Gamma$.

  \tagitem{prvFalse}{$\lfalse\in\Gamma$ olsaydı
    $\Gamma\Proves[\Sigma]\lfalse$ olur ve $\Gamma$
    $\Sigma$-tutarsız olurdu.}{}

  \tagitem{prvTrue}{$\Gamma\Proves[\Sigma]\ltrue$ olduğundan türetim
    altında kapalılık, yani \olref{prop:ccs-closed} durumu gereğince
    $\ltrue\in\Gamma$.}{}

\item $\lnot!A\in\Gamma$ ise tutarlılık gereğince $!A\notin\Gamma$;
    tersine $!A\notin\Gamma$ ise $\Gamma$ tam $\Sigma$-tutarlı olduğundan
    $\lnot!A\in\Gamma$.

  \tagitem{prvAnd}{\iftag{probAnd}{Alıştırma.}{$!A\land !B\in\Gamma$
      olsun. $(!A\land !B)\lif !A$ bir totoloji örneği olduğundan türetim
      altında kapalılık, yani \olref{prop:ccs-closed} durumu gereğince
      $!A\in\Gamma$. $!B\in\Gamma$ için de benzer biçimde. Öte yandan hem
      $!A\in\Gamma$ hem $!B\in\Gamma$ olsun. Bu durumda
      $!A\lif(!B\lif(!A\land !B))$ bir totoloji örneği olduğundan türetim
      altında kapalılık $(!A\land !B)\in\Gamma$ olmasını gerektirir.}}{}

  \tagitem{prvOr}{\iftag{probOr}{Alıştırma.}{$!A\lor !B\in\Gamma$,
      $!A\notin\Gamma$ ve $!B\notin\Gamma$ olsun. $\Gamma$ tam
      $\Sigma$-tutarlı olduğundan $\lnot!A\in\Gamma$ ve
      $\lnot!B\in\Gamma$. $\lnot!A\lif(\lnot!B\lif\lnot(!A\lor !B))$
      bir totoloji örneği olduğundan $\lnot(!A\lor !B)\in\Gamma$ olurdu.
      Bu, $\Gamma$'nın $\Sigma$-tutarsız olduğu anlamına gelir; çelişki.}}{}

  \tagitem{prvIf}{\iftag{probIf}{Alıştırma.}{$!A\lif !B\in\Gamma$ ve
      $!A\in\Gamma$ olsun; o zaman $\Gamma\Proves[\Sigma]!B$ ve türetim
      altında kapalılık gereğince $!B\in\Gamma$. Tersine
      $!A\lif !B\notin\Gamma$ ise $\Gamma$ tam $\Sigma$-tutarlı olduğundan
      $\lnot(!A\lif !B)\in\Gamma$. $\lnot(!A\lif !B)\lif !A$ bir totoloji
      örneği olduğundan $!A\in\Gamma$; $\lnot(!A\lif !B)\lif\lnot !B$ bir
      totoloji örneği olduğundan $\lnot!B\in\Gamma$. Tutarlılık gereğince
      $!B\notin\Gamma$.}}{}

  \tagitem{prvIff}{\iftag{probIff}{Alıştırma.}{$!A\liff !B\in\Gamma$
      olsun. $!A\in\Gamma$ ise $(!A\liff !B)\lif(!A\lif !B)$ bir totoloji
      örneği olduğundan $!B\in\Gamma$. $!B\in\Gamma$ ise benzer biçimde
      $!A\in\Gamma$. Böylece ya ikisi de $\Gamma$ içindedir ya da hiçbiri
      değildir.

      Tersine $!A\liff !B\notin\Gamma$ olsun. $\Gamma$ tam
      $\Sigma$-tutarlı olduğundan $\lnot(!A\liff !B)\in\Gamma$.
      $\lnot(!A\liff !B)\lif(!A\lif\lnot !B)$ bir totoloji örneği
      olduğundan, $!A\in\Gamma$ ise $\lnot!B\in\Gamma$ ve tutarlılık
      gereğince $!B\notin\Gamma$. Benzer biçimde $!B\in\Gamma$ ise
      $!A\notin\Gamma$. Dolayısıyla ne $!A\in\Gamma$ ve $!B\in\Gamma$ ne
      de $!A\notin\Gamma$ ve $!B\notin\Gamma$ olur.}}{}
  \end{enumerate}
\end{proof}

\begin{probtag}{probAnd,probOr,probIf,probIff}
  \olref[nml][com][ccs]{prop:ccs-properties} önermesinin ispatını tamamlayın.
\end{probtag}

